\section{Detector con Alarma Sonora}
\begin{itemize}
  \item[(c)] El alumno usará un detector con alarma sonora, relacionará la señal audible con el índice de exposición correspondiente y reconocerá las ventajas de su uso continuo. (0:30 hr)
\end{itemize}

\section{Práctica No. 3 -- Manejo y uso del detector con alarma sonora}
\begin{itemize}
  \item[(1 hr)]
\end{itemize}

\subsection{Instrumentos de monitoreo}
\textbf{Objetivo General.} -- El alumno manejará correctamente al menos uno de los siguientes instrumentos de monitoreo con cámara de ionización, con detector Geiger o con detector de centelleo. Interpretará correctamente las lecturas tomando en consideración las limitaciones de cada uno. Podrá verificar la calibración del instrumento usando una fuente de referencia.

\begin{enumerate}
  \item Monitor con detector Geiger
    \begin{itemize}
      \item Lo manejará y explicará en qué casos es conveniente su uso. Describirá sus ventajas y desventajas frente a una cámara de ionización. Describirá el fenómeno de saturación y el uso del factor de calibración.
    \end{itemize}
\end{enumerate}

\section{Práctica No. 4 -- Manejo y uso de un monitor con detector Geiger}
\begin{itemize}
  \item[(1 hr)]
\end{itemize}

\section{Práctica No. 5 -- Localización de fuentes con Geiger-Müller}
\begin{itemize}
  \item[(1:30 hr)]
\end{itemize}

\section{Filosofía de la Protección Radiológica}
\begin{itemize}
  \item[(2 hr)]
\end{itemize}

\subsection{Concepto de Riesgo-Beneficio. Concepto ALARA.}
El alumno podrá percatarse de las diferencias entre el riesgo y el beneficio, en diversos tipos de actividades y, en especial, en la radiografía industrial. Explicará en forma simplificada el concepto ALARA.

\subsection{Sistema de Limitación de Dosis}
Memorizará el valor del límite de equivalente de dosis para trabajadores ocupacionalmente expuestos y los valores de rapidez de exposición en mrem/hr.
